\section{简易逻辑}
\subsection{从集合到命题}
命题最最基础的定义是“可以判断真假的陈述句。”\index{命题}，在一定的范围内可以判定真假是判断是否是命题的核心方法。
一个误区是所有的命题都可以写成若p则q的形式，这是不对的。实际上命题的核心形式是$x \in \mathbb{X} $,很显然元素属于集合与否是容易判断的。
但是我们见到的大多数命题都经过了各种修饰与复合，使得看起来和集合已经没有什么关系了。
我们最常见的两种命题就是若p则q和带有$\forall,\exists $的命题。（当然两者结合在一起也可以）
\clearpage
\subsection{四种命题*}

\clearpage
\subsection{充分条件，必要条件，充要条件}
\begin{definition}
    若“若p，则q”是真命题，则记作$p \Rightarrow q$，称p是q的充分条件\index{充分条件}，q是p的必要条件\index{必要条件}。
反之，若“若p，则q”是假命题，则记作$p \not\Rightarrow  q$，称p不是q的充分条件，q不是p的必要条件。
\end{definition}

思考：充分一词很容易理解，那么必要一词如何理解？即q不成立，p一定不成立。请用一些具体的命题来验证一下这个说法。
\begin{definition}
    若满足$p \Rightarrow q$且$q \Rightarrow p$，则称p与q互为充分必要条件，简称充要条件。
记作
\[p \Leftrightarrow q\]
\end{definition}

\begin{exercise}
    请写出“p是q的充分不必要条件”和“p是q的必要不充分条件”的等价命题。
\end{exercise}


与集合的关系
\\若$p \Rightarrow q$，其中$p$是$x \in \mathbb{P} $，$q$是$x \in \mathbb{Q} $,
则$\mathbb{P} $与$\mathbb{Q} $的关系是？（$\mathbb{P}\subseteq\mathbb{Q} $）简记为：在小范围里一定在大范围里

例8
\begin{example}
    已知集合$\mathbb{P} =\{x|a-4<x<a+4\}$，$\mathbb{Q} =\{x|1<x<3\}$,“$x \in \mathbb{P} $”是“$x \in \mathbb{Q} $”
的必要条件，则实数a的取值范围是？
\end{example}

\begin{jie}
    “$x \in \mathbb{P} $”是“$x \in \mathbb{Q} $”的必要条件，即“$x \in \mathbb{Q} $”是“$x \in \mathbb{P} $”
的充分条件，根据上文可得$\mathbb{Q} \subseteq \mathbb{P} $.列出不等式组

\[  
    \left\{
    \begin{lgathered}
    a-4\leq 1\\
    3\leq a+4\\
    \end{lgathered}
    \right.
\]
最后解得$a\in [-1,5]$
\end{jie}
\begin{exercise}
    判断一下说法的正误：
    \begin{enumerate}
        \item $a=b$是$ac=bc$的充分不必要条件；
        \item $\frac{1}{a}>\frac{1}{b}$是$a<b$的既不充分也不必要条件；
        \item $a\neq 0$是$ab \neq 0$的必要不充分条件；
        \item $a>b>0$是$a^n>b^n(n\in \mathbb{N}\text{且}n\leq2)$的充要条件。
    \end{enumerate}
\end{exercise}
\begin{exercise}
    已知条件p:$|x-1|\leq 2$，条件q:$x>a$，且满足q是p的必要不充分条件，则a取值为？
\end{exercise}

\subsection{全称量词、存在量词与命题的否定}
\begin{definition}
    在逻辑中，诸如“所有的”，“任意一个”的短语叫做全称量词\index{全称量词}，用记号$\forall$表示；
同理，诸如“存在”“至少一个”的短语叫做存在量词（特称量词）\index{存在量词}，用记号$\exists$表示。记作
\[\forall x\in M,P(x).\exists x\in M,P(x).\]
其含义也可以看做是
\[\forall x\in M,x\in \{x|P(x)\}.\exists x\in M,x\in \{x|P(x)\}.\]
\end{definition}
翻译条件
\begin{example}
    $\forall x\in A,\exists y\in B : y=x$的等价条件是？
\end{example}

\begin{jie}
    \[A\subseteq B\]
\end{jie}


“你有多大脚，我有多大鞋。”


练习：借助函数图像，使用最大值与最小值的语言翻译下列叙述。
\\$\forall x_1 \in D_1,\forall x_2 \in D_2,f(x_1)>g(x_2).\Leftrightarrow$
\\$\exists x_1 \in D_1,\exists x_2 \in D_2,f(x_1)>g(x_2).\Leftrightarrow$
\\$\forall x_1 \in D_1,\exists x_2 \in D_2,f(x_1)>g(x_2).\Leftrightarrow$
\\$\exists x_1 \in D_1,\forall x_2 \in D_2,f(x_1)>g(x_2).\Leftrightarrow$
\\$\forall x \in D$,使得$f(x)>g(x)$恒成立$\Leftrightarrow$
\\$\exists x \in D$,使得$f(x)>g(x)$可以成立$\Leftrightarrow$
\\$\forall x_1 \in D_1,\exists x_2 \in D_2,f(x_1)=g(x_2).\Leftrightarrow$
\\$\exists x_1 \in D_1,\exists x_2 \in D_2,f(x_1)=g(x_2).\Leftrightarrow$


首先应当引入一个概念——排中律，即一个命题与它的否定应该有且仅有一个成立，
有时候我们也成命题的否定为原命题的“非命题”，这写都意味着原命题与命题的否定必然是一对一错的关系。
而我们主要研究的就是带有量词的。
\begin{theorem}
    随着$\neg$划过整个命题，命题中的“$\exists$”和$“\forall$”互换，最终停留在表达式前面。
\end{theorem}

\begin{example}
    \[\forall x\in M,P(x).\exists x\in M,P(x).\]
对应的否定是
\[\exists x\in M,\neg P(x).\forall x\in M,\neg P(x).\]
\end{example}


多个量词自然也不难，应用上一节的练习验证这个结论。

思考：全程量词和存在量词的否定的结果是彼此互换的逻辑是什么？是否可以帮助我们证明或否定一些含有量词的命题？

前面量词部分的转换是容易的，重点还是后面从$P(X)$到$\neg P(x)$的做法。
\\在前文我们提到过，逻辑用语的最本质形式是$x \in \mathbb{X} $，
而我们要做的就是把$\in$改为$\notin$即可。相应的，也就是将$=$改为$\neq$,$>$改为$\leq$

思考：若p，则q的形式的命题如何改成否定形式呢？

\clearpage

